\documentclass[12pt,notitlepage]{report}
% The shared notes preamble lives outside this repo.
\input{/Users/aryankulkarni/Documents/y3-notes/preamble.tex}

\title{Week 5: residuals and diagnostics for the injury model}
\author{Aryan Kulkarni}
\date{MTH6134 group project}

\begin{document}
\maketitle

\paragraph{Findings.}
The injury model of week 4 passes its checks overall, and no player drives
it. One weakness shows: among young players on high training loads the
product term predicts too few injuries. A variant in which load matters only
above age 25 removes it. The linear model for minutes was checked in week 2
and stands.

\section{Residuals for a GLM}

In a GLM the variance changes with the mean, so a raw residual
$y_i - \hat\mu_i$ has to be scaled before residuals can be compared. The two
standard choices are
\[
  r^P_i = \frac{y_i - \hat\mu_i}{\sqrt{V(\hat\mu_i)}},
  \qquad
  r^D_i = \operatorname{sign}(y_i - \hat\mu_i)\sqrt{d_i},
\]
the Pearson and deviance residuals, where $d_i$ is observation $i$'s
contribution to the deviance. Dividing either by $\sqrt{1 - h_{ii}}$
standardises it, with $h_{ii}$ the leverage from the final weighted
least-squares fit.

For a binary response these fail as plotting tools. Given $\hat\mu_i$, a
residual takes one of two values, so a plot against fitted values is two
curves whatever the model. Two repairs are used here.

\paragraph{Binned residuals.}
Sort the players by fitted probability, cut them into 20 bins and average the
raw residuals in each. A bin mean has standard error
$\sqrt{\sum \hat\pi_i(1 - \hat\pi_i)}/n_{\text{bin}}$ under the model, so
about 19 of 20 bins should lie within two standard errors of zero.

\paragraph{Randomised quantile residuals.}
Let $F$ be the fitted distribution function of $Y_i$. Draw
$u_i \dist U\{F(y_i - 1), F(y_i)\}$ and set $r^Q_i = \Phi^{-1}(u_i)$. If the
model is right these are exactly standard normal, so the plots of week 2
apply again. The function is \mintinline{r}|qresid()| in
\mintinline{r}|common.R|.

\begin{minted}{r}
source("../common.R")
players <- players |> mutate(age_c = age - 25, load_c = training_load - 3)
fit_main <- glm(
  injured ~ age_c + load_c + prev_injuries,
  family = binomial,
  data = players
)
fit_inj <- update(fit_main, . ~ . + age_c:load_c)

binned <- function(fit, bins = 20) {
  tibble(fitted = fitted(fit), residual = fit$y - fitted(fit)) |>
    mutate(bin = ntile(fitted, bins)) |>
    group_by(bin) |>
    summarise(
      mean_fitted = mean(fitted),
      mean_residual = mean(residual),
      bound = 2 * sqrt(sum(fitted * (1 - fitted))) / n()
    )
}
bins <- binned(fit_inj)
q <- qresid(fit_inj)

base_fig("residuals", mfrow = c(1, 3), height = 2.3)
plot(
  predict(fit_inj),
  resid(fit_inj),
  col = pt,
  xlab = "Linear predictor",
  ylab = "Deviance residuals"
)
plot(
  bins$mean_fitted,
  bins$mean_residual,
  col = blue,
  ylim = range(-bins$bound, bins$bound, bins$mean_residual),
  xlab = "Mean fitted probability",
  ylab = "Mean residual in bin"
)
abline(h = 0)
lines(bins$mean_fitted, bins$bound, lty = 2)
lines(bins$mean_fitted, -bins$bound, lty = 2)
qqnorm(q, main = "", col = pt)
qqline(q, lwd = 2)
dev.off()

out(
  c(
    bins_outside = sum(abs(bins$mean_residual) > bins$bound),
    largest_abs_quantile_residual = max(abs(q))
  ),
  "residual-numbers"
)
\end{minted}

\begin{figure}[H]
  \centering
  \incfig{residuals}
  \caption{Left: deviance residuals form two bands and say nothing. Centre:
    binned residuals with two-standard-error bounds. Right: QQ plot of
  quantile residuals.}
  \label{fig:residuals}
\end{figure}

\inputminted{text}{output/residual-numbers.txt}

One of the 20 bins lies outside its bounds, which is the number expected by
chance, and the bin means show no trend across the range of fitted risk. The
quantile residuals follow the normal line.

\section{Can the checks find a missing term?}

A diagnostic earns trust by catching a known fault. The main-effects model
omits the age by load interaction, so its residuals should be positive where
age and load are both high.

\begin{minted}{r}
cells <- players |>
  mutate(
    age_band = cut(
      age,
      c(16, 22, 25, 28, 36),
      labels = c("17-22", "23-25", "26-28", "29-36")
    ),
    load = ifelse(training_load >= 4, "4-5", "1-3"),
    main = injured - fitted(fit_main),
    final = injured - fitted(fit_inj),
    v = fitted(fit_inj) * (1 - fitted(fit_inj))
  ) |>
  group_by(load, age_band) |>
  summarise(
    n = n(),
    main = mean(main),
    final = mean(final),
    two_se = 2 * sqrt(sum(v)) / n(),
    .groups = "drop"
  )
out(
  as.data.frame(cells) |> mutate(across(main:two_se, \(x) round(x, 3))),
  "cells"
)
\end{minted}

\inputminted{text}{output/cells.txt}

The main-effects model is caught. At high load its mean residual climbs with
age from $-0.046$ to $+0.163$, and the last cell lies outside its bound of
0.139: among players aged 29 and over on loads 4 to 5, the injury rate is 16
percentage points above what that model fits. With the interaction included
all eight cells are inside their bounds.

\section{Leverage and influence}

\begin{minted}{r}
cook <- cooks.distance(fit_inj)
cook_cut <- 8 / (nobs(fit_inj) - 2 * length(coef(fit_inj)))
flagged <- which(cook > cook_cut)
out(
  c(
    threshold = cook_cut,
    flagged = length(flagged),
    flagged_and_injured = sum(players$injured[flagged]),
    largest_cook = max(cook),
    largest_fitted_risk_among_flagged = max(fitted(fit_inj)[flagged])
  ),
  "influence"
)

fit_drop <- update(fit_inj, data = players[-flagged, ])
out(
  round(cbind(all = coef(fit_inj), without = coef(fit_drop)), 3),
  "sensitivity"
)
\end{minted}

\inputminted{text}{output/influence.txt}
\inputminted{text}{output/sensitivity.txt}

Fourteen players exceed the threshold of 0.013, none by much (the largest
distance is 0.035). Thirteen are injured players to whom the model gave a low
risk, and one is an uninjured player with a fitted risk of 0.77. In a binary
model the influential observations are the surprises. Without them every
coefficient grows in size, the interaction from 0.111 to 0.139, so these
players pull the estimates towards zero and no conclusion rests on them. They
stay in.

\section{Calibration}

The Hosmer-Lemeshow statistic groups players by fitted risk and compares
observed with expected injuries in each group, as a Pearson $\chi^2$ on
(groups $-$ 2) degrees of freedom.

\begin{minted}{r}
hosmer <- function(fit, groups) {
  p <- fitted(fit)
  group <- cut(p, unique(quantile(p, 0:groups / groups)), include.lowest = TRUE)
  observed <- tapply(fit$y, group, sum)
  expected <- tapply(p, group, sum)
  size <- tapply(p, group, length)
  x2 <- sum((observed - expected)^2 / (expected * (1 - expected / size)))
  c(
    groups = groups,
    x2 = x2,
    p = pchisq(x2, length(observed) - 2, lower.tail = FALSE)
  )
}
out(
  round(t(sapply(c(5, 6, 8, 10, 12, 15), hosmer, fit = fit_inj)), 3),
  "hosmer"
)

fit_alt <- glm(
  injured ~ age_c + prev_injuries + I(pmax(age_c, 0) * load_c),
  family = binomial,
  data = players
)
young <- players$age <= 21 & players$training_load >= 4
compare <- rbind(
  product = c(
    AIC = AIC(fit_inj),
    hosmer(fit_inj, 10)["p"],
    expected = sum(fitted(fit_inj)[young])
  ),
  above_25 = c(
    AIC = AIC(fit_alt),
    hosmer(fit_alt, 10)["p"],
    expected = sum(fitted(fit_alt)[young])
  )
)
out(round(compare, 3), "compare")
out(
  c(players = sum(young), injured = sum(players$injured[young])),
  "young-high-load"
)
\end{minted}

\inputminted{text}{output/hosmer.txt}
\inputminted{text}{output/young-high-load.txt}
\inputminted{text}{output/compare.txt}

The verdict depends on the number of groups: $p$ is 0.71 with 5 groups,
0.003 with 10 and 0.32 with 12. A statistic this unstable cannot settle the
question of fit, but the 10-group result points somewhere specific. Of the 37
players aged 21 or under on loads 4 to 5, 5 were injured where the model
expects 2.6. This is the falling part of the curves in week 4: the product
term lowers the fitted risk of young players as load rises, and the data show
no such fall.

The second model lets load act only above age 25, through the term
$\max(\text{age} - 25, 0) \times \text{load}$. It has one parameter fewer, a
lower AIC (532.4 against 535.3), expects 4.1 injuries in that group and
passes the 10-group test ($p = 0.36$). It was chosen after seeing the misfit,
so its support is overstated. It tells the same story in a more plausible
form: load carries no extra risk up to 25 and a growing one after.

\section{Summary}

\begin{itemize}
  \item Residual plots for binary data need binning or randomisation. With
    them the injury model shows no systematic misfit across the range of
    fitted risk.
  \item No single player or small group of players drives the conclusions.
  \item The product form of the interaction understates risk for young
    players on high loads. The variant with a load effect from age 25 fits
    at least as well and is a candidate for the report.
\end{itemize}

\end{document}
